\chapter{Reference Implementation}

This document is associated with a reference implementation available at:

\begin{center}
\url{https://github.com/your-repository/globalplatform-open}
\end{center}

The current code base is organized as a small workspace with three main
Rust crates:

\begin{itemize}
\item \texttt{core}, a reusable embedded library crate that provides the
  runtime, target dispatch, serial, semihosting, flash and crypto
  services;
\item \texttt{gp\_kernel}, the production firmware crate that owns the
  current APDU main loop and application dispatch;
\item \texttt{core\_test}, a dedicated embedded test firmware used to
  exercise the shared runtime under QEMU and report machine-readable
  verdicts through semihosting.
\end{itemize}

This separation avoids mixing production entry points, reusable
services, and test-only firmware logic in a single crate. It also makes
it possible to run host-side \texttt{cargo test} orchestration against
real embedded binaries executed under QEMU, rather than against a hosted
simulation of the embedded code.

\section{APDU Transport over T=0}
\label{sec:impl-apdu-t0}
% --- APDU T=0 transport and mapping ---
% This section describes how the abstract APDU model is mapped onto the T=0 protocol.
% It should explain:
% - APDU cases (Case 1–4)
% - byte-level exchange (procedure bytes, SW1/SW2 timing)
% - handling of Lc/Le
% - fragmentation strategy for LOAD
% - interaction with secure messaging (wrapping before transport)

The current reference implementation establishes a small but explicit
\texttt{T=0}-style exchange model before introducing full command
semantics. At boot, the runtime emits an ATR in direct convention
(\texttt{TS = 3B}), announces a single interface byte through
\texttt{T0 = 1C}, publishes the default serial-rate byte
\texttt{TA1 = 11}, carries twelve historical bytes, and then enters the
main APDU loop.

As specified in Section~\ref{sec:apdu-serial-adaptation}, the current boards carry
this exchange over a conventional serial link; they do not reproduce the
ISO/IEC~7816 contact-card electrical interface. The byte transport remains
behind \texttt{TransportLayer} so that an ISO/IEC~7816-capable target can
provide the reference physical link without changing the \texttt{T=0} state
machine or the APDU layers above it.

Those historical bytes are encoded as ISO compact-TLV data:
\begin{verbatim}
80 56 09 C1 DE 5E xx yy 73 80 00 00
\end{verbatim}
The first byte, \texttt{80}, is the category indicator that makes the
remaining historical bytes a compact-TLV sequence. The object
\texttt{56 09 C1 DE 5E xx yy} uses compact tag~5, length~6, as issuer data:
\texttt{09 C1 DE 5E} is the Oxide SE marker, \texttt{xx} is the
OS version byte, and \texttt{yy} is a build-capability bitmap.
In the current build \texttt{xx = 09}, meaning version 0.9 beta. In the
\texttt{yy} bitmap, bit~7
is set when the root Security Domain is
not \texttt{NullSecurityDomain}; bits~6,~5 and~4 advertise compiled and
enabled SCP11a, SCP11b and SCP11c support; bits~3 and~2 are reserved;
and bits~1--0 encode SCP03 support (\texttt{00}: none, \texttt{10}:
S8, \texttt{11}: S16).

The object \texttt{73 80 00 00} is the ISO Card Capabilities compact
TLV, tag~7 and length~3. Oxide SE deliberately advertises only the
minimal capability it really supports here: selection by full DF name,
which matches SELECT by AID. It does not advertise an ISO file system,
record or binary EF commands, extended APDUs, or ISO logical-channel
management.

Each loop iteration reads only the five command-header bytes
(\texttt{CLA}, \texttt{INS}, \texttt{P1}, \texttt{P2}, and the fifth
length byte) and constructs one transport-owned APDU session. The base
\texttt{T=0} layer does not attempt to decide on its own whether that
fifth byte is
\texttt{Lc} or \texttt{Le}. Instead, command logic receives an
APDU-facing object and explicitly requests incoming or outgoing
behavior. This choice follows the same idea as Java Card:
\texttt{T=0} transport first captures the command header, while the
command logic decides whether it is about to consume command data or to
prepare response data.

\begin{figure}[ht]
  \centering
  \begin{tikzpicture}[
    stage/.style={
      draw,
      rounded corners=2pt,
      minimum width=3.1cm,
      minimum height=0.9cm,
      align=center,
      font=\footnotesize
    },
    arrow/.style={->, thick, >=Latex}
  ]
    \node[stage] (atr) {Boot\\send ATR};
    \node[stage, right=1.0cm of atr] (hdr) {Read only the\\5-byte APDU header};
    \node[stage, right=1.0cm of hdr] (proc) {Build \texttt{Apdu}\\call \texttt{process\_apdu()}};
    \node[stage, below left=1.25cm and 0.55cm of proc] (direct) {Outgoing only\\send bytes now\\then append status};
    \node[stage, below right=1.25cm and 0.55cm of proc] (deferred) {Incoming consumed and\\outgoing prepared\\store bytes and return \texttt{61xx}};
    \node[stage, below=1.25cm of deferred] (getresp) {\texttt{GET RESPONSE}\\send next chunk\\then status or \texttt{61xx}};

    \draw[arrow] (atr) -- (hdr);
    \draw[arrow] (hdr) -- (proc);
    \draw[arrow] (proc) -- (direct);
    \draw[arrow] (proc) -- (deferred);
    \draw[arrow] (deferred) -- (getresp);
  \end{tikzpicture}
  \caption{Current \texttt{T=0} APDU exchange flow in the reference implementation.}
  \label{fig:impl-apdu-exchange}
\end{figure}

The transport-owned APDU session keeps the five header bytes in
dedicated scalar fields and exposes a separate fixed data buffer of
\texttt{256} bytes
for command or response data. This matches the current short-APDU
scope, avoids dynamic allocation on the transport path, and stays close
to the Java Card idea of a single APDU work buffer. The base
\texttt{T=0} manager retains ownership of actual byte emission
decisions, while command
handlers only observe an APDU-facing surface that can request
``incoming'', ``outgoing'', and ``outgoing length'' transitions.

\subsection{APDU Implementation Surfaces}
\label{sec:impl-apdu-surfaces}

At this stage, the implementation deliberately separates APDU handling
into six levels, each with a different responsibility.

\begin{itemize}
\item \textbf{Byte transport.} The \texttt{TransportLayer} trait owns
  only byte emission and byte reception. In the current reference
  implementation it is instantiated either as a serial backend or as a
  scripted in-kernel simulator.
\item \textbf{Base \texttt{T=0} manager.} The
  \texttt{T0ApduManager<T>} object is built on top of a
  \texttt{TransportLayer}. It owns the current \texttt{T=0} exchange
  state: the five header bytes, the transport buffer, the deferred
  response buffer, and the logic for \texttt{61xx},
  \texttt{GET RESPONSE}, and final \texttt{SW1/SW2} emission.
\item \textbf{Composable APDU layers.} The \texttt{ApduLayer} trait is
  the protocol-layer boundary above \texttt{T=0}. It exposes command
  reception and completion through \texttt{ApduCommand} and
  \texttt{ApduCompletion}, so higher protocol layers do not need to
  know the concrete transport-owned APDU representation.
\item \textbf{Kernel APDU view.} The kernel exposes a narrower
  \texttt{Apdu} view over that transport state. This is the object used
  by kernel-side command handlers such as the current
  \texttt{kernel\_main\_app}.
\item \textbf{Shared Rustlet ABI buffer.} When a selected Rustlet is
  called, the kernel stages the current command into a fixed-layout
  \texttt{RustletCtx}. This ABI object contains the header fields, the
  current status word, direction metadata, and one fixed payload area.
\item \textbf{Common APDU trait.} Both the kernel-side \texttt{Apdu}
  view and the Rustlet-side \texttt{RustletCtx} implement the common
  \texttt{SEApdu} trait. This is the normalization point of the current
  implementation: kernel handlers, Rustlet handlers, and the
  kernel-to-Rustlet bridge expose the same conceptual APDU operations.
\end{itemize}

This layering is important. \texttt{SEApdu} is not the transport object
itself. \texttt{ApduLayer} is not the card-side processing interface
either. \texttt{SEApdu} is the card-side processing interface of ``the
APDU currently being handled''. In other words:
\begin{itemize}
\item \texttt{TransportLayer} owns bytes on the wire;
\item \texttt{T0ApduManager} owns the \texttt{T=0} state machine;
\item \texttt{ApduLayer} owns protocol-layer composition above
  \texttt{T=0};
\item \texttt{SEApdu} owns the command-processing view seen by command
  logic.
\end{itemize}

The present implementation therefore keeps transport semantics in the
kernel while still giving both kernel-side code and Rustlets a common
APDU vocabulary. This is the main reason why the current
\texttt{kernel\_main\_app} and selected-Rustlet dispatch path can be
written against the same card-side APDU concepts while the byte-level
\texttt{T=0} policy remains centralized below them.

The current runtime stack instantiated by the kernel is therefore:
\begin{center}
\texttt{TransportLayer $\rightarrow$ T0ApduManager $\rightarrow$ ApduLayer wrappers $\rightarrow$ dispatcher $\rightarrow$ SEApdu}
\end{center}

At the time of writing, the wrapper chain contains a pass-through layer,
the generic secure-channel layer, and a small tracing layer used for scripted
debug runs. The secure-channel layer owns the APDU boundary for SCP03 and
SCP11 establishment. It also owns protected-command unwrapping and
protected-response wrapping.

\subsection{Implementation Impact of Timing Constraints}

Section~\ref{sec:apdu-t0-timing} explained the transport-level timing
constraint in general terms. In the present implementation, this leads
to a very concrete coding rule: the main loop must keep the
transmit/receive transition points adjacent in the control flow.

In particular, the implementation intentionally treats the following
boundaries as hot:

\begin{itemize}
\item after the ATR has been emitted, the code must branch immediately
  to the command-reception loop;
\item after the five-byte header has been read, the code must hand
  control immediately to the APDU handler so that an incoming payload
  can be received without unrelated work being interposed;
\item after \texttt{SW1/SW2} has been emitted, the code must fall back
  immediately to the next command reception;
\item after a deferred-response \texttt{61xx} status, the code must be
  ready to receive \texttt{GET RESPONSE} without adding unrelated work
  on the path.
\end{itemize}

This is why the current \texttt{gp\_kernel} main loop is deliberately
minimal and places comments directly at these transition points in the
source code. The goal is not to claim formal timing conformance from a
single software structure, but to make the hot path explicit and to
avoid accidental regressions caused by debug output, bookkeeping, or
other apparently harmless additions at the wrong place.

The current APDU-facing helpers implement the following contract:

\begin{itemize}
\item \texttt{set\_incoming\_and\_Receive()} reads exactly
  \texttt{Lc} bytes from the serial link into the APDU data buffer and
  marks the APDU as having consumed incoming data;
\item \texttt{set\_outgoing()} marks the APDU as an outgoing exchange
  and returns the current \texttt{Le} interpretation;
\item \texttt{set\_outgoing\_length(length)} declares how many outgoing
  bytes are prepared and constrains the writable outgoing slice to that
  length.
\end{itemize}

These names are the intended public command-processing vocabulary. The
kernel uses idiomatic Rust internally where useful, while the conceptual
APDU interface is intentionally aligned with Java Card.
The Rustlet runtime also provides \texttt{send\_bytes()} and
\texttt{send\_bytes\_long()} as convenience operations layered on top of the
same APDU buffer and outgoing-length declaration.

This last point is the main \texttt{T=0} design constraint of the
current implementation. A command is allowed to read incoming data and
to prepare outgoing data in the same logical transaction, but the
response bytes are never emitted directly by the handler itself.
Instead, the base \texttt{T=0} manager checks two facts after command
processing returns:

\begin{itemize}
\item whether the APDU object has prepared outgoing data;
\item whether the same APDU has already consumed incoming data.
\end{itemize}

If outgoing data exists and no incoming payload was consumed, the base
\texttt{T=0} manager sends the prepared bytes immediately and then appends
\texttt{SW1/SW2}. If outgoing data exists and incoming data was
consumed, the manager copies the response bytes into its own pending
response buffer, returns a status word of the form \texttt{61xx}, and
waits for a subsequent \texttt{GET RESPONSE} command. The deferred
buffer is therefore owned by the \texttt{T=0} manager rather than by the
\texttt{Apdu} object itself.

The current \texttt{GET RESPONSE} recognition rule is intentionally
small:

\begin{itemize}
\item \texttt{INS = C0};
\item \texttt{P1 = 00};
\item \texttt{P2 = 00}.
\end{itemize}

When such a command is received, the \texttt{T=0} manager handles it
itself rather than forwarding it to command logic. The
pending response buffer then emits up to the requested number of bytes
and either:

\begin{itemize}
\item returns the final completion status when the buffered response has
  been fully drained;
\item or returns another \texttt{61xx} if additional response bytes
  remain to be fetched.
\end{itemize}

In the current short-buffer implementation, a \texttt{GET RESPONSE}
command with \texttt{Le = 00} is interpreted as a request for
\texttt{256} bytes, matching the conventional \texttt{T=0} short-length
encoding.

The expected status values at this stage are therefore:

\begin{itemize}
\item \texttt{9000}: normal success, used when a command completes
  without deferred response data, or after the last successful
  \texttt{GET RESPONSE} chunk has been sent;
\item \texttt{61xx}: response bytes are pending and must be retrieved by
  \texttt{GET RESPONSE}; \texttt{SW2} carries the currently remaining
  byte count, with \texttt{00} representing \texttt{256} remaining
  bytes in the current fixed-buffer implementation;
\item \texttt{6985}: conditions not satisfied; used when a
  \texttt{GET RESPONSE} is received without any pending response, or
  when another command is received while deferred response bytes are
  still waiting to be fetched;
\item \texttt{6D00}: instruction not supported; currently the default
  result returned by \texttt{process\_apdu()} for all unimplemented
  commands.
\end{itemize}

This yields three distinct completion paths:

\begin{verbatim}
// Direct outgoing path
let le = apdu.set_outgoing();
apdu.set_outgoing_length(response_len);
let buffer = apdu.buffer_mut();
buffer[..response_len].copy_from_slice(response_data);
// T0ApduManager then appends SW1/SW2
\end{verbatim}

\begin{verbatim}
// Deferred T=0 path
apdu.set_incoming_and_Receive();
let le = apdu.set_outgoing();
apdu.set_outgoing_length(response_len);
let buffer = apdu.buffer_mut();
buffer[..response_len].copy_from_slice(response_data);
// T0ApduManager stores the response, returns 61xx,
// then waits for GET RESPONSE
\end{verbatim}

\begin{verbatim}
// Wrong Le for a pure outgoing command
let le = apdu.set_outgoing();
apdu.set_outgoing_length(actual_len);
// T0ApduManager compares requested Le and actual_len
// and may return 6Cxx instead of the payload
\end{verbatim}

The detailed correspondence between these APDU-facing operations and the
underlying \texttt{T=0} byte stream is summarized in
Appendix~\ref{app:apdu-over-t0}.

\section{Application Lifecycle Management}
\label{sec:impl-lifecycle}

Chapter~4 defines the lifecycle contract retained by Oxide SE. The reference
kernel implements that contract at the object-registry boundary: every managed
object carries a kind-specific lifecycle state, and APDU handlers consult that
state before allowing the operation family that would make the object effective.

The implemented state families are deliberately small:
\begin{itemize}
\item packages are \texttt{Loaded} or \texttt{Locked}; only
      \texttt{Loaded} packages may be used by \texttt{INSTALL [for install]};
\item ordinary instances are \texttt{Selectable} or \texttt{Locked}; only
      \texttt{Selectable} instances may be selected;
\item Security Domain instances are \texttt{Selectable} or \texttt{Locked};
      only \texttt{Selectable} Security Domains may be selected or used to
      establish a secure channel;
\item key objects are \texttt{Active} or \texttt{Locked}; only
      \texttt{Active} keys may be consumed by secure-channel establishment.
\end{itemize}

Deletion is represented by absence from the registry. Failed or interrupted
deployment is represented by the absence of a newly committed registry entry.
This avoids publishing half-created packages or instances as stable managed
objects.

\subsection{LOAD Command Handling}
\label{sec:impl-load}

The normative \texttt{INSTALL [for load]} and \texttt{LOAD} APDU shapes are
described in Section~\ref{sec:gp-install-command}. The kernel implements them
as one volatile dynamic-load transaction. \texttt{INSTALL [for load]} records
the selected authority Security Domain, the target Security Domain, the package
AID, the announced size, and the expected hash. It does not publish a package
object in the registry.

Each following \texttt{LOAD} command must target the same authority and arrive
with the next expected block number in \texttt{P2}. The final block is marked
by \texttt{P1}. The kernel rejects out-of-order blocks, size overflow beyond
the announced length, authority mismatch, and final hash mismatch. Any APDU
that is not part of the active load sequence cancels the volatile context.

Only the final valid \texttt{LOAD} block commits the package as a registry
object in state \texttt{Loaded}. Because the registry is mutated only at this
final point, a reset or power loss during the load sequence falls back to the
previous committed registry view.

\subsection{INSTALL Command Handling}
\label{sec:impl-install}

The normative command family is described in
Section~\ref{sec:gp-install-command}. This subsection only describes the
reference-kernel path for the implemented \texttt{INSTALL [for install]} case:
binding an already known package to a new selectable instance, optionally
creating a Security-Domain instance, and initializing the Rustlet state that
will later be serialized in the registry.

The kernel accepts the \texttt{INSTALL [for install]} shape identified by
\texttt{P1 = 0C} and \texttt{P2 = 00}. The command data is parsed as the
GlobalPlatform length-value sequence:
\begin{itemize}
\item package AID;
\item applet or module AID;
\item application instance AID;
\item privilege bytes;
\item install parameters.
\end{itemize}

Empty package, applet, or instance AIDs are rejected, and the privilege field
is bounded to the three-byte administrative form used by the Security-Domain
model described in Chapter~4. Malformed LV data is rejected before any registry
mutation or Rustlet call is attempted.

The management authority is resolved before installation. A clear
\texttt{INSTALL} command is interpreted under the root Security Domain
authority, and is accepted only if that backend explicitly permits out-of-band
management. A protected \texttt{INSTALL} command is interpreted under the
Security Domain bound to the authenticated secure-channel session, after the
unwrap path described in Section~\ref{sec:impl-scp-ext}. In both cases, the
kernel asks the authority to approve the operation family before it mutates the
registry. These checks cover Security-Domain and registry-plane access,
management of the target package/applet/instance triple, making the target
selectable, and the profile-specific installation policy hook.

If the requested object is a kernel-native Security Domain instance, the
kernel does not enter a Rustlet image. It creates the Security-Domain registry
object directly, attaches it to the authority Security Domain through
\texttt{parent\_sd\_aid}, and stores the encoded administrative state derived
from the install privilege bytes. This is the path used for
\texttt{NullSecurityDomain} and \texttt{KernelSecurityDomain} instances.

Otherwise, installation targets a Rustlet package already present in the
registry. The kernel resolves the package/applet pair in the active
Security-Domain view, activates the selected Rustlet image as the new instance,
and fills the shared Rustlet context with a compact LV payload containing the
same package AID, applet AID, instance AID, privilege bytes, and install
parameters. The Rustlet runtime receives that payload through its
\texttt{install()} handler, as described in
Section~\ref{sec:impl-rustlet-runtime}. The implementation deliberately avoids
building a second transport APDU for this call; the installation payload is
encoded directly into the shared Rustlet ABI buffer.

After a successful Rustlet \texttt{install()} return, the kernel saves the
instance state into the global object registry under the authority Security
Domain. Ordinary Rustlet instances are stored as instance objects with
serialized state. Rustlet-backed Security Domains additionally expose the
Security-Domain dispatch vtable described in Section~\ref{sec:impl-sd-profiles}
and keep their kernel-enforced privilege floor in the Security-Domain registry
object. A failed install call, a failed authorization check, or a failed state
save leaves no successfully selectable instance behind.

At this layer, lifecycle handling is intentionally direct: successful package
loading creates a \texttt{Loaded} package, successful installation creates a
\texttt{Selectable} instance or Security Domain, and authorized
\texttt{SET STATUS} commands lock or unlock those objects. \texttt{SELECT}
itself only changes the current execution context.

\subsection{DELETE Command Handling}
\label{sec:impl-delete}
% --- DELETE command handling ---
% This subsection describes how the DELETE command is implemented.
% It should cover:
% - validation of AIDs to delete
% - dependency checking (applications vs packages)
% - prevention of deletion of active or selected applications
% - registry update and consistency enforcement
% - memory and resource reclamation
% - rollback or failure handling

\tbc

\subsection{Lifecycle State Machine}
\label{sec:impl-lifecycle-sm}

The kernel stores lifecycle state inside the polymorphic object-registry
payload, not in a parallel table. The state bits are persisted with registry
entries, so a locked package, instance, Security Domain, or key remains locked
after reboot when persistent storage is available.

The current transition causes are:
\begin{itemize}
\item successful dynamic load commits a package as \texttt{Loaded};
\item successful \texttt{INSTALL [for install]} commits a
      \texttt{Selectable} ordinary instance or Security Domain;
\item successful \texttt{PUT KEY} commits a key object as \texttt{Active};
\item \texttt{DELETE} removes the addressed registry object after authority
      and dependency checks;
\item \texttt{SET STATUS} toggles packages, ordinary instances, and Security
      Domains between their usable state and \texttt{Locked}.
\end{itemize}

The guards are checked where the protected operation starts. Package state is
checked during package lookup for \texttt{INSTALL [for install]}; instance and
Security Domain selectability is checked during \texttt{SELECT};
Security-Domain selectability is checked before secure-channel establishment;
key state is checked while resolving SCP03 key material.

\texttt{SET STATUS} is handled as a management command. The APDU layer resolves
the authority exactly as for other management operations: a protected command is
interpreted under the Security Domain bound to the secure channel, while a
clear command is interpreted under the root Security Domain. The target object
must be reachable from that authority. Clear \texttt{SET STATUS} is accepted
only when the authority explicitly permits it; this is true for
\texttt{NullSecurityDomain} development images and false for kernel-native and
Rustlet-backed production Security Domains. Rustlet-backed Security Domains
receive a \texttt{set\_status()} authorization hook through the same dispatch
bridge used by \texttt{INSTALL}, \texttt{DELETE}, \texttt{PUT KEY}, and
\texttt{STORE DATA}.

\subsection{GET STATUS Implementation}
% --- GET STATUS implementation ---
\label{sec:impl-get-status}

The implementation parses only the modern GlobalPlatform TLV form described
in Section~\ref{sec:gp-get-data-discovery}. The mandatory \texttt{4F} search
criterion is canonicalized into a bounded AID value; the optional \texttt{5C}
tag list is structurally validated. Unsupported object families or response
formats fail before registry traversal.

Traversal follows stable registry slot order. For every live slot, the kernel
checks the requested family, AID prefix, lifecycle state and visibility from
the active Security Domain authority. Package, ordinary-instance and
Security-Domain objects are encoded directly as \texttt{E3} records in the
APDU response buffer. Key and generic-data objects are never candidates.

If the next matching record does not fit, the response ends with
\texttt{63 10}. The continuation cursor contains only:

\begin{itemize}
\item the authority AID;
\item the canonical query;
\item the index of the next registry slot.
\end{itemize}

The next command must repeat that query with \texttt{P2 = 03}. A different
query, a different authority, or any intervening non-\texttt{GET STATUS}
command invalidates the cursor. This prevents continuation under a changed
security context and avoids retaining a copy of any registry object.

The encoder is allocation-free. It first computes the exact short BER-TLV
record length, then writes directly into the remaining APDU payload. Thus
pagination is determined by the transport capacity rather than by an
intermediate response array.

% --- Cryptographic primitive layer ---
% This section describes:
% - minimal crypto API exposed by core
% - opaque AES key representation
% - in-place buffer-oriented operations
% - target-specific hardware hooks
% - software fallback strategy
\section{Cryptographic Primitive Backend}
\label{sec:impl-crypto-backend}

Before discussing key storage and SCP03 session handling, the reference
implementation needs a small cryptographic primitive layer. This layer is not
meant to expose a large application-facing crypto API. Its role is narrower:
to provide the primitive operations required by SCP03 and by platform key
management, while keeping the implementation compatible with constrained
embedded targets and with future hardware-assisted ports.

The current design therefore adopts three explicit choices.

First, the public API is buffer-oriented and keeps allocation out of the
cryptographic path. Command processing and secure-channel wrapping already
operate on preallocated APDU buffers, so the crypto layer should follow the
same model.

Second, AES keys are represented through an opaque \texttt{AesKey} type. The
caller can construct such a key from raw bytes, but does not manipulate the
internal representation directly. This leaves room for later target-specific
formats, key-slot indirections, or hardware key handles without changing the
higher-level SCP03 code.

Third, the backend selection is intentionally asymmetric:
\begin{itemize}
\item a target-specific hardware implementation should be preferred whenever a
  serious port exposes constant-time AES support or a dedicated cryptographic
  peripheral;
\item a software implementation must remain available as a fallback, both for
  boards without accelerators and for bring-up under emulation.
\end{itemize}

For the current project stage, the minimal API retained in \texttt{core} is:
\begin{itemize}
\item \texttt{fill\_random(buf)} for entropy or challenge generation;
\item \texttt{aes\_cbc\_encrypt\_in\_place(key, iv, buf)} for block-aligned
  in-place encryption;
\item \texttt{aes\_cbc\_decrypt\_in\_place(key, iv, buf)} for block-aligned
  in-place decryption;
\item \texttt{aes\_cmac(key, input, out)} for CMAC computation;
\item \texttt{scp03\_kdf(...)} for the SCP03 session-derivation path, with the
  derivation constant, context, and output buffer provided explicitly by the
  caller.
\end{itemize}

The in-place form of the CBC helpers is deliberate. In a constrained runtime,
it is highly desirable that input and output may be the same buffer. This keeps
the interface close to the APDU transport model and avoids needless scratch
buffers when wrapping or unwrapping secure messages.

The current reference implementation also assumes that the present QEMU targets
do not expose a cryptographic accelerator usable by the project at this stage.
As a consequence, the portable fallback path is the relevant one for current
development and test environments.

The runtime distinguishes explicitly between the board itself and the
execution environment. The build flow injects:
\begin{itemize}
\item \texttt{OXIDE\_SE\_BOARD} to select the supported target board;
\item \texttt{OXIDE\_SE\_EXECUTION\_ENV} to distinguish a hardware-oriented build
  from a QEMU-oriented one.
\end{itemize}

This second dimension matters because some services depend on the actual
execution context rather than on the nominal board identity alone. The same
\texttt{STM32F405} board support, for instance, should use the on-chip RNG on
real hardware but must not assume that the corresponding QEMU model exposes
that peripheral.

Among available open implementations, the RustCrypto AES crate is a credible
reference for that fallback path \cite{rustcrypto-aes}. Its public
documentation states that its implementations are designed to execute in
constant time, either through hardware intrinsics when available or through a
portable bitsliced backend otherwise. This does not eliminate the need for
review, integration discipline, and platform-specific testing, but it provides
a reasonable software baseline for the present project stage.

The random-generation path is treated more conservatively than the AES path.
For entropy generation, the project does not currently accept a purely
deterministic software fallback with no external entropy source. The NIST
\texttt{SP 800-90} family explicitly separates the role of entropy sources from
that of deterministic generators: a DRBG construction is expected to be fed by
an entropy source rather than to replace it \cite{nist80090b,nist80090c}. For
that reason, the default behavior of \texttt{fill\_random()} remains an
explicit error whenever no credible entropy source has been wired for the
current target.

For the supported STM32 hardware targets, the intended production path is to
use the on-chip RNG peripheral directly. Both STM32F405 and STM32L475 expose a
dedicated RNG block with a control register, a status register, and a data
register. The usage pattern is straightforward:
\begin{itemize}
\item initialize the peripheral once during the crypto-backend startup path;
\item enable the appropriate clock source for the target;
\item enable the RNG peripheral;
\item wait until the data-ready flag indicates that a fresh 32-bit word is
  available;
\item read the data register;
\item reject or reinitialize on hardware entropy error flags.
\end{itemize}

In practice, the generation latency documented by ST is on the order of a few
dozen peripheral clock cycles per 32-bit word, which is small enough that the
reference design does not currently require any software ``diversification'' or
PRNG expansion layer on top of the hardware output. A simple implementation may
therefore block until the status bit becomes ready and then return the newly
read word. A more optimized implementation may keep one prefetched word in
flight in order to overlap the wait for the next sample, but this remains an
optimization of the same basic hardware-driven design.

The QEMU situation is different. The currently used STM32 board models do not
provide the corresponding RNG device \cite{qemu-stm32}. The present
implementation therefore does not try to emulate a fake TRNG in software.
Instead, QEMU-oriented builds use a distinct development path selected through
\texttt{OXIDE\_SE\_EXECUTION\_ENV=qemu}: \texttt{fill\_random()} is serviced through
semihosting file access to the host \texttt{/dev/urandom}. This keeps the
development path explicit and separate from the production hardware path, while
avoiding the introduction of an ad hoc software-only generator inside the
guest.

This QEMU-only entropy bridge should not be confused with a production entropy
source. It is a bring-up facility that depends on host cooperation and on the
availability of semihosting. The default policy for unsupported environments
still remains an explicit error rather than a silently invented random stream.

More generally, the existence of a QEMU execution mode must not be read as a
security claim about the emulated target. The current QEMU path is a prototype
and validation aid. It is intentionally useful for debugging, but precisely for
that reason it relies on mechanisms that are inappropriate as part of a secure
deployment story. In particular, semihosting gives the guest controlled access
to host services and is explicitly documented by QEMU as a development-oriented
facility rather than as an isolation-preserving boundary
\cite{qemu-security,qemu-semihosting}. The QEMU runtime is therefore a
convenience environment for bring-up and test, not a substitute for the
security properties expected from the real hardware platform.

This mirrors another already known limitation of the current QEMU bring-up,
namely the lack of usable flash support in the emulated environment.

From an architectural point of view, this section should be read as a lower
layer than the next two ones:
\begin{itemize}
\item the key-management layer decides how static SCP03 keys are stored,
  versioned and updated;
\item the secure-channel layer consumes those keys and invokes the primitive
  backend to derive \texttt{S-ENC}, \texttt{S-MAC}, and \texttt{S-RMAC}, then
  to compute \texttt{C-MAC}, \texttt{R-MAC}, and protected payloads.
\end{itemize}

In a mature hardware port, the same high-level API should therefore dispatch to
target hooks that exploit the actual silicon capabilities, while preserving the
software fallback for development, unsupported boards, and test environments.

% --- Key management and secure storage ---
% This section describes:
% - internal representation of keys (key sets, versions, identifiers)
% - secure storage mechanisms
% - key update policies (PUT KEY)
% - diversification and derivation
% - protection against leakage
\section{Key Management and Secure Storage}
\label{sec:impl-key-management}

The current implementation uses the global registry for resident packages,
application instances, Security Domain instances, and cryptographic key
objects attached to those Security Domains. At the time of writing, the first
implemented key-object family is SCP03 static communication keys.

This means that the registry stores four broad classes of objects:
\begin{itemize}
\item packages carrying executable binary code;
\item ordinary application instances carrying serialized state;
\item Security Domain instances and their administrative state;
\item typed key objects attached to a Security Domain instance, currently
      SCP03 static keys.
\end{itemize}

For these ancillary key objects, the ownership rule is the same as for other
managed entities: each object is attached to one owning Security Domain
instance through \texttt{parent\_sd\_aid}. Key metadata is held in dedicated
registry fields, while the payload contains the raw key bytes consumed by the
secure-channel implementation.

In the current subset, one SCP03 key object records:
\begin{itemize}
\item a key-version number;
\item a key identifier / keyset identifier;
\item a usage selector (\texttt{ENC} or \texttt{MAC});
\item a key state, currently \texttt{Active} or \texttt{Locked};
\item one AES-128 static key value.
\end{itemize}

The current \texttt{PUT KEY} implementation therefore behaves as follows:
\begin{itemize}
\item the active Security Domain authorizes the command through the normal
      management boundary;
\item the kernel decodes the GlobalPlatform reference controls: the low seven
      bits of \texttt{P1} carry the key version, the high bit of \texttt{P1}
      carries the last-command flag, the low seven bits of \texttt{P2} carry
      the first key identifier, and the high bit of \texttt{P2} indicates
      whether the command contains several keys;
\item the kernel parses the payload entries
      \texttt{(usage, key\_len, key\_bytes)};
\item each accepted entry is stored as one serialized registry object attached
      to the active Security Domain instance;
\item when several payload entries are present, their effective key
      identifiers increment from the first identifier decoded from
      \texttt{P2};
\item storing an already existing
      \texttt{(parent\_sd\_aid, key\_version, key\_id, usage)} tuple rotates
      the key material by replacing the registry object atomically;
\item later \texttt{INITIALIZE UPDATE} resolves the requested
      \texttt{(key\_version, key\_id)} pair by reading back the corresponding
      \texttt{ENC} and \texttt{MAC} objects from the registry and rebuilding
      the \texttt{StaticKeys} structure consumed by the pure SCP03 engine.
\end{itemize}

\texttt{Locked} key objects remain present in the registry for policy and
audit-oriented evolution, but they are ignored by SCP03 session derivation.
Deletion helpers operate on the same owner-qualified tuple, so deleting one
usage or keyset member cannot remove a sibling Security Domain's key.

This storage rule is shared by both secure backends currently implemented in
Oxide SE:
\begin{itemize}
\item kernel-native Security Domains resolve those registry-backed key objects
      directly inside the kernel;
\item Rustlet-backed Security Domains reached through
      \texttt{RustletSecurityDomainProxy} load the same key objects through
      runtime syscalls and derive their SCP03 sessions in user-land.
\end{itemize}

Two consequences are important:
\begin{itemize}
\item SCP03 static key material is stored as Security Domain-owned registry
      objects;
\item identical \texttt{(key\_version, key\_id, usage)} tuples may exist in
      several Security Domain subtrees, because the corresponding key objects
      are qualified by \texttt{parent\_sd\_aid}.
\end{itemize}

This organization keeps an important implementation invariant intact:
cryptographic material is not hidden in ad hoc side tables unrelated to the
registry model. The registry remains the common persistence and authority
surface for applications, Security Domains, and the ancillary objects they
own.

For the root kernel-native Security Domain, the bootstrap path already inserts
the static SCP03 key objects explicitly declared under \texttt{root.keys} in
the build manifest. Supplementary Security Domains use the equivalent
\texttt{security\_domains.keys} declaration. The very first
\texttt{INITIALIZE UPDATE} therefore consumes registry-backed key material
rather than a side-loaded session secret, and a manifest without keys receives
no implicit development fallback. The same registry model is also used by the
Rustlet-backed Security Domain path; its difference is not the storage model,
but where session derivation is executed.

The remaining work in this area is not the existence of registry-backed keys
themselves, but the richer GlobalPlatform policy surface around them:
additional usages, rotation semantics, diversification, production
provisioning, and the final policy for first-key provisioning of supplementary
Security Domain instances.


\section{Secure Channel Implementation}
\label{sec:impl-scp}

Oxide SE implements secure channels as a kernel-mediated protocol stack whose
active Security Domain backend may be kernel-native or Rustlet-backed. The
current implementation is intentionally split into four layers:

\begin{itemize}
\item protocol-specific pure cryptographic engines
      (\texttt{core::scp03}, \texttt{core::scp11}) with no dependency on APDU
      routing, registry state, or Security Domain policy;
\item \texttt{core::gp\_sm}, which centralizes the bounded payload budgets
      imposed by Amendment~D command and response MAC trailers;
\item the active Security Domain backend, which owns policy decisions, static
      keys, session lifecycle, and secure-messaging state;
\item \texttt{SecureChannelLayer}, which owns the APDU boundary: command
      interception, command unwrapping, response wrapping, and management
      command authorization for the supported SCP profiles.
\end{itemize}

This split is the main implementation invariant. Cryptographic transforms are
not scattered through the APDU dispatcher, and the cryptographic core does not
know how an AID is installed, selected, or stored in the registry.

The active backend may take three forms:
\begin{itemize}
\item a pseudo-Security-Domain path (\texttt{NullSecurityDomain}) for bring-up
      and test scenarios;
\item a kernel-native Security Domain path (\texttt{KernelSecurityDomain});
\item a Rustlet-backed Security Domain path reached through
      \texttt{RustletSecurityDomainProxy}.
\end{itemize}

In every case, the secure-channel session is associated with one Security
Domain instance and one logical channel. The current implementation still
provides only one effective channel, but the authority boundary is already
framed in those terms rather than as a single process-global SCP state.

\subsection{Security Domain Profiles}
\label{sec:impl-sd-profiles}
\label{sec:impl-scp-profiles}

The default management selector is configured by the predeployment manifest.
The root Security Domain backend is declared under \texttt{[root.package]},
while the secure-channel protocols compiled into the image are declared under
\texttt{[secure\_channel]}. Three Security Domain backends are currently
accepted:

\begin{itemize}
\item \texttt{NullSecurityDomain}, a no-security/debug authority that
      authorizes management operations and does not establish SCP sessions;
\item \texttt{KernelSecurityDomain}, a kernel-native Security Domain whose
      supported SCP profiles are selected by \texttt{protocols},
      \texttt{scp03\_profile}, and \texttt{scp11\_profiles};
\item \texttt{RustletSecurityDomainProxy}, a kernel-side proxy that delegates
      Security-Domain policy and SCP state to one selected Rustlet Security
      Domain instance while preserving kernel-owned base privileges.
\end{itemize}

The S8 and S16 variants are explicit secure-channel profiles rather than
distinct administrative authorities. A manifest that enables SCP03 therefore
declares \texttt{scp03\_profile = "S8"} or
\texttt{scp03\_profile = "S16"}. A manifest that enables only SCP11 omits the
SCP03 profile entirely and declares its supported SCP11 variants through
\texttt{scp11\_profiles}, for example \texttt{["A", "C"]}. The internal engine
is shared, and one \texttt{KernelSecurityDomain} implementation exposes the
same \texttt{SecurityDomainManagement} and
\texttt{SecurityDomainSecureChannel} traits regardless of the selected SCP
combination.

The proxy variant keeps those same kernel-side traits and APDU boundaries, but
forwards the dedicated Security-Domain operations through \texttt{sddispatch}
to the currently selected Rustlet Security Domain instance. The kernel still
owns APDU parsing, registry mutation, non-escalation checks, and active-context
tracking; the Rustlet Security Domain refines policy and performs its own SCP
cryptography through the runtime APIs.

\subsection{Selected Secure-Channel Security Levels}
\label{sec:impl-scp-selected-levels}

The reference implementation deliberately selects a strict subset of the
normative security-level matrices described in Chapter~3. This is a build and
interoperability profile, not a redefinition of the GlobalPlatform values.

For SCP03, \texttt{EXTERNAL AUTHENTICATE} accepts only
\texttt{P1 = 00}, \texttt{01}, or \texttt{03}: authenticated session without
secure messaging, C-MAC, or C-MAC with C-ENC. The response-protection levels
\texttt{11}, \texttt{13}, and \texttt{33} are rejected explicitly with
\texttt{6985}. Consequently, responses in the selected SCP03 profile remain
plain; the optional \texttt{BEGIN R-MAC SESSION} and
\texttt{END R-MAC SESSION} commands are profiled out and return
\texttt{6D00}.

For SCP11a, SCP11b, and SCP11c, the CRT accepts only key-usage qualifier
\texttt{95 01 3C}, selecting C-MAC, C-ENC, R-MAC, and R-ENC. The GP-defined
\texttt{34} alternative is rejected for all three variants; the SCP11c-only
\texttt{1C} and \texttt{74} alternatives are rejected as well. SCP11 does not
define the explicit R-MAC session companion commands: response protection is
part of the security level established by the CRT.

The authenticated-peer class is also part of the selected profile and is
applied before any Security Domain hook:
\begin{itemize}
\item SCP03 and owner-authenticated SCP11a may carry the implemented
      management command families;
\item SCP11b authenticates the card, not the OCE, and is consequently limited
      to read-only \texttt{GET DATA} in this implementation profile;
\item owner-authenticated SCP11c permits the implemented management families
      except \texttt{PUT KEY}, \texttt{SET STATUS}, and deletion of key
      objects, as required by the SCP11c restrictions;
\item SCP11c \texttt{ANY\_AUTH} permits \texttt{GET DATA} only. Mutating
      operations requiring the \texttt{BF20} authorization object are
      explicitly profiled out.
\end{itemize}
This matrix is a non-overridable kernel gate. The selected Security Domain's
hierarchy, privilege checks, and management hooks are then evaluated for every
accepted command and may only narrow the decision.

The implementation also supports a supplementary-protocol fallback. A
secure-channel protocol selected in the build manifest always has priority and
uses the kernel matrix above. If no selected kernel profile recognizes an
establishment command, the active Rustlet Security Domain may explicitly claim
the header and process the raw establishment data. The resulting session is
recorded as Rustlet-owned without assigning it a fictitious SCP identifier.
Protocol-specific authorization is then owned by that Security Domain, but
registry ancestry, privilege non-escalation, lifecycle checks, APDU bounds, and
management-family hooks remain enforced by the kernel. This extension point
therefore permits supplementary protocols without turning the registry into
Rustlet-controlled memory.

These choices are enforced before a session becomes authenticated. Host-based
tests exhaustively reject every other byte value, while QEMU scenarios exercise
the accepted levels and representative GP-defined but profiled-out values at
the APDU boundary.

S8 uses 8-byte host/card challenges, 8-byte authentication cryptograms, and
8-byte C-MAC/R-MAC values. S16 uses 16-byte challenges, 16-byte cryptograms,
and 16-byte MAC values. The APDU layer therefore never assumes a fixed
protected-command trailer shape. It asks the active Security Domain for the
current MAC length and session policy before separating the raw payload from
its trailing MAC.

Separately from these kernel-native profiles, Oxide SE also supports a
Rustlet-backed Security Domain path. In that case, the kernel still owns APDU
interception and routing, but the selected Security Domain instance performs
secure-channel operations through the Rustlet runtime and the generated
\texttt{sddispatch} entry.

\subsection{Pure SCP03 Engine}
\label{sec:impl-scp-core}

The pure engine in \texttt{core::scp03} defines:

\begin{itemize}
\item \texttt{Scp03Profile}, with profile-dependent lengths for S8 and S16;
\item \texttt{StaticKeys}, currently AES-128 ENC and MAC test keys;
\item \texttt{SessionKeys}, containing S-ENC, S-MAC, and S-RMAC;
\item \texttt{SessionMaterial}, produced by one
      \texttt{INITIALIZE UPDATE};
\item \texttt{SessionState}, which owns MAC chains and encryption counters for
      secure messaging;
\item \texttt{SecurityLevel}, which represents no secure messaging, SCP03
      command-only protection, and the full bidirectional protection selected
      for SCP11.
\end{itemize}

The engine derives session material from the static keys, host challenge and
card challenge using the SCP03 KDF. It then computes card and host
cryptograms, verifies the host cryptogram during authentication, validates
C-MAC on protected commands, decrypts C-ENC command data, and provides the
R-MAC/R-ENC primitives used by the common SCP11 secure-messaging path. The
selected SCP03 levels do not activate response protection; the exact selected
matrix is stated in Section~\ref{sec:impl-scp-selected-levels}.

The current implementation uses AES-CMAC and AES-CBC style encryption with
GlobalPlatform padding rules. SCP03 itself remains a symmetric protocol here:
asymmetric authentication, post-quantum key establishment, or card certificate
handling would be higher-level platform services, not requirements of the
current SCP03 core.

\subsection{APDU Layer Integration}
\label{sec:impl-scp-apdu-layer}

\texttt{SecureChannelLayer} sits above the T=0 APDU manager. It sees clear APDUs
after T=0 procedure-byte handling, but before Rustlet dispatch or management
mutation. Its responsibilities are:

\begin{itemize}
\item intercept \texttt{INITIALIZE UPDATE} and call the active Security
      Domain's secure-channel entry point;
\item intercept \texttt{EXTERNAL AUTHENTICATE}, verify the host cryptogram,
      and open the requested security level;
\item intercept \texttt{PERFORM SECURITY OPERATION} and
      \texttt{MUTUAL AUTHENTICATE} for SCP11-capable backends;
\item reject protected commands if no authenticated secure-channel session
      exists;
\item split protected command data into its raw payload and trailing
      profile-dependent C-MAC;
\item ask the Security Domain to unwrap protected command data before command
      dispatch;
\item ask the Security Domain to protect response data, append the R-MAC when
      required, and retain the actual status word before transmission.
\end{itemize}

Management commands such as \texttt{INSTALL [for install]} are still parsed and
executed by the kernel. The Security Domain authorizes them. With
\texttt{NullSecurityDomain}, this authorization is deliberately open for debug
and bootstrap scenarios. With \texttt{KernelSecurityDomain}, mutating
management commands require an
authenticated encrypted SCP03 session. The kernel-native Security Domain path
keeps a small per-instance administrative state keyed by instance
AID. At this stage that state serializes the privilege bitfield carried by
\texttt{INSTALL}, seeds the implicit root instance with the bootstrap privilege
set, and lets the kernel preserve kernel-owned administrative privilege state
without relying on a Rustlet serializer.

More generally, Security Domain instances are modeled as objects of the global
registry. Each such object may carry a \texttt{parent\_sd\_aid} and a
serialized administrative state. For a Rustlet-backed Security Domain, that
state is split between kernel-owned authority data and Rustlet-owned persistent
functional state. For \texttt{NullSecurityDomain} and kernel-native profiles,
the persistent administrative state may reduce to privilege bits. Active
secure-channel state remains volatile in every backend.

\subsection{Persistent Security-Domain State and Volatile Sessions}
\label{sec:impl-scp-state-lifetime}

The lifetime of Security-Domain configuration is distinct from the lifetime of
an authenticated secure-channel session. Registry-backed persistent state may
contain lifecycle, privileges, application policy, certificate references, or
other long-lived configuration. It shall not contain derived session keys,
expected cryptograms, receipts, command or response MAC chains, encryption
counters, or partially staged handshake buffers.

A Rustlet Security Domain expresses this boundary in its serialized Rust type.
Its secure-channel field is marked \texttt{\#[serde(skip)]} and implements
\texttt{Default}. The runtime deserializes the object when the Security Domain
FAE is loaded. The resulting Rust value remains alive in RAM while that
Security Domain is loaded, so consecutive establishment and protected APDUs
observe the same live session. The active Rustlet Security Domain remains
resident in a dedicated firmware slot while an ordinary selected Rustlet
occupies the application slot. The \texttt{unwrap}, application-dispatch, and
\texttt{wrap} calls therefore alternate between two live instances rather than
unloading and restoring the Security Domain through the registry. After each
successful \texttt{sddispatch} call, the runtime serializes only the persistent
view and publishes it through the registry. It does not reload persistent state
before every APDU.

Unloading, rebooting, or later reloading the Security Domain reconstructs the
skipped session field at its inactive default. An earlier Postcard encoding
that still carries trailing session bytes may be decoded for migration, but
those bytes are ignored and never reactivate the old channel.

Explicit \texttt{reset\_secure\_channel()} remains mandatory. A new
establishment exchange, a terminating clear command, or another session-ending
condition may occur while the Security Domain remains loaded. Reset therefore
overwrites all volatile secrets, chains, counters, cryptograms, receipts, and
staged handshake data immediately; session termination never relies on a
future deselection. Persistent mutations are still committed after successful
calls rather than being deferred until \texttt{DESELECT}, because power may be
lost before such a command occurs.

\subsection{Validation Strategy}
\label{sec:impl-scp-validation}

Validation is deliberately split into two levels. Pure cryptographic regression
tests validate deterministic SCP03 transformations for S8 and S16. QEMU
integration tests validate that those transformations are correctly connected
to APDU parsing, Security Domain policy, registry mutation, and Rustlet
execution.

The detailed vector files, their OpenSCDP/GlobalPlatformPro and Samsung
OpenSCP-Java provenance, and the exact test flows are documented in
Appendix~\ref{app:scp03-vectors}. The
important implementation point here is that the kernel does not treat S8/S16
as test-only constants: profile-dependent lengths are part of the active
Security Domain state and are consumed by the APDU secure-channel layer.

\subsection{Remaining Certification Gap}
\label{sec:impl-scp-certification-gap}

The implementation is structured for auditability, but it is not a claim of
GlobalPlatform certification. Appendix~\ref{app:scp03-vectors} documents the
current independent regression evidence. The remaining work includes:

\begin{itemize}
\item defining a secure production provisioning policy for the initial keys
      currently declared by development manifests;
\item validating ICV, chaining, counter and error behavior against official
      certification material if it becomes available to the project;
\item completing key-set lifecycle policy beyond the current registry-backed
      \texttt{PUT KEY} path, including production provisioning, rotation
      governance and diversification rules.
\end{itemize}

\section{Registry and Data Management}
\label{sec:impl-registry}
% --- Internal registry and data access ---
% This section describes:
% - internal representation of registry (AIDs, domains, privileges)
% - GET DATA / STORE DATA handling
% - persistence model
% - consistency guarantees

\subsection{GlobalPlatform Discovery Objects}
\label{sec:impl-gp-discovery}

The APDU layer resolves \texttt{GET DATA 0066} and \texttt{0067} through the
Security Domain that is authoritative for the command. Clear discovery uses
the root Security Domain; protected discovery uses the Security Domain bound
to the active secure channel.

\texttt{NullSecurityDomain} emits a complete \texttt{66/73} recognition
structure but no \texttt{64} SCP declaration, and reports
\texttt{GET DATA 0067} as absent with \texttt{6A88}.
\texttt{KernelSecurityDomain} constructs both objects from the effective build
configuration: only compiled and enabled SCP03 and SCP11 profiles are
advertised. SCP11 always advertises S16, while SCP03 reports its configured S8
or S16 option.

A Rustlet-backed Security Domain answers through its \texttt{get\_data} hook.
The runtime default constructs the same standards-compliant objects from the
SCP03 S8/S16 and SCP11a/b/c support methods. A Rustlet may override
\texttt{get\_data} for richer domain-specific data. If it advertises no
protocol, the proxy does not substitute the kernel's capabilities: discovery
describes the selected Security Domain, not the firmware in general.

All discovery objects are encoded directly in the APDU response buffer by a
bounded BER-TLV writer. No registry object, key bytes, certificate, or
intermediate response array is exposed or retained.

The firmware keeps one kernel-owned object registry. This registry is the
implementation counterpart of the GlobalPlatform object model described in
Chapter~4: it is the single catalogue for packages, application instances,
Security Domain instances, and ancillary objects such as secure-channel keys.
Security Domains do not own private registries. Instead, every registry object
is attached to an administrative Security Domain through
\texttt{parent\_sd\_aid}.

Each registry entry has a small common header:
\begin{itemize}
\item \texttt{object\_aid}, the AID by which the object is addressed;
\item \texttt{parent\_sd\_aid}, the Security Domain instance under whose
  administrative authority the object lives;
\item \texttt{kind}, which distinguishes the concrete object family.
\end{itemize}

The object kind determines the payload shape:
\begin{itemize}
\item a \textbf{package} object uses its \texttt{object\_aid} as the
  executable-load-file AID and carries the FAE binary image;
\item an \textbf{application instance} records the package AID from which it
  was instantiated and carries the serialized application state;
\item a \textbf{Security Domain instance} records the package or backend that
  implements it, the three-byte administrative privilege field, and its
  serialized administrative state;
\item a \textbf{key} object records key metadata such as version, identifier,
  usage, type and length, together with the raw key material.
\end{itemize}

The root Security Domain is not kept outside this structure. It is created by
the predeployment manifest as a real Security Domain instance with no parent,
using the backend declared under \texttt{[root.package]} and the instance AID
and installation bytes declared under \texttt{[root.instance]}. A development
image may use \texttt{NullSecurityDomain}; a protected image may use
\texttt{KernelSecurityDomain} or a Rustlet-backed Security Domain. In all
cases, the root follows the same registry shape as supplementary Security
Domains.

Registry lookup is authority-aware. A raw global lookup exists only inside the
kernel for consistency checks and bootstrap. Management operations resolve AIDs
through the active Security Domain view: the selected or secure-channel-bound
Security Domain sees itself, objects attached to it, and objects attached to
descendant Security Domains. Objects outside that subtree are treated as absent
for that administrative operation. This is the implementation form of the
visibility rule described in Section~\ref{sec:gp-object-relationships}.

Privilege delegation is checked before a Security Domain instance is created or
updated. The three-byte privilege field is kernel-owned authority data: a
Rustlet-backed Security Domain may receive the installation payload and keep
additional serialized policy state, but it cannot broaden the base
GlobalPlatform privileges that the kernel recorded for the instance. A child
Security Domain must therefore have a privilege set no broader than what its
parent is allowed to delegate.

Persistence is deliberately expressed at the registry-object level. Packages,
instances, Security Domains, and keys are the units that must eventually be
written atomically to non-volatile storage. Until the persistent backend is
implemented, the predeployment manifest rebuilds the initial registry at boot.
Once non-volatile storage is available, the same registry-object model remains
the unit of recovery: a management operation either installs, updates, or
deletes a complete registry object consistently, or it leaves the previous
registry state intact.

\section{Logical Channel and Context Management}
\label{sec:impl-channel}
% --- Logical channels and execution context ---
% This section describes:
% - channel allocation (MANAGE CHANNEL)
% - mapping channel → selected application
% - SELECT command handling
% - context switching model

The current kernel implements only the minimal selected-context model
required to bootstrap an embedded administrative application. In
GlobalPlatform terms, this is deliberately only the beginning of the
logical-channel story: there is one effective channel, one selected
context slot in the kernel, and no \texttt{MANAGE CHANNEL}
implementation yet.

At this stage, the selected context is not a high-level Rust object but
an explicit kernel-owned record containing:
\begin{itemize}
\item the application-global-pointer value returned by the loaded FAE
  image;
\item an opaque application state pointer;
\item a pointer to a C-style vtable containing the application entry
  points.
\end{itemize}

This representation is intentional. The kernel image and the loaded FAE
image are two distinct position-independent binaries, each with its own
global-pointer context. A native Rust \texttt{\&dyn Trait} would not be
an appropriate boundary here, because the call path must explicitly
switch to the application's relocation context before entering its code
and must restore the kernel context on return.

\subsection{Selected Context and Embedded Administrative Application Activation}
\label{sec:impl-selected-context}

The selected application is drawn from Rustlets packaged as independent FAE
images and declared by the predeployment manifest. Each registry package and
instance uses three GlobalPlatform-style identifiers:
\begin{itemize}
\item a \textbf{package AID}, corresponding to the executable load file;
\item an \textbf{applet AID}, corresponding to the installable module;
\item an \textbf{instance AID}, corresponding to the selectable live
  instance.
\end{itemize}

Test manifests may assign the same value to these three AIDs for compact
scenarios. They remain distinct identifiers in the registry and follow the
GlobalPlatform identification model.

The \texttt{SELECT} command remains kernel-owned. The transport loop
detects \texttt{INS = A4} before generic dispatch, requests the command
payload, interprets it as an instance AID, and compares it against the
kernel registry. If no entry matches, the kernel returns \texttt{6A82}
(file not found). If a matching entry exists, the kernel activates the
corresponding embedded FAE and records that instance as the current
selected context.

This activation follows the XiPFS FAE execution model rather than a
kernel-specific reinvention of the startup runtime. The kernel enters
the embedded image at offset \texttt{0}, thereby using the existing
XiPFS startup code unchanged. The startup path relocates the FAE image
into the application RAM arena and yields the information required to
call the application \texttt{start} entry point.

The kernel-side activation path therefore has two distinct steps:
\begin{enumerate}
\item run the FAE loader to relocate the image and obtain the
  application entry point together with the application global-pointer
  state;
\item call the application's \texttt{start(buffer)}
  function and validate the returned selected-application descriptor.
\end{enumerate}

The \texttt{buffer} currently contains the shared ABI buffer used to
carry the ABI version, APDU header fields, status word and payload. The
descriptor returned by \texttt{start()} contains:
\begin{itemize}
\item the ABI version expected by the application;
\item the current instance AID as seen from the Rustlet runtime;
\item an opaque application state pointer;
\item a vtable with two entry points:
  \texttt{install()} and \texttt{process\_apdu()}.
\end{itemize}

On the application side, the current ABI is intentionally minimalist.
An application declares one resident state type and one installation
function, then exports them through:

\begin{verbatim}
struct SecurityDomain;
declare_rustlet!(SecurityDomain, 1024usize, install);
\end{verbatim}

The \texttt{declare\_rustlet!} macro hides the internal runtime
boilerplate in \texttt{rustlet\_runtime}. The
application author therefore only needs to provide:
\begin{itemize}
\item the concrete state type later used by
  \texttt{process\_apdu()};
\item an \texttt{install()} function that constructs this state;
\item the \texttt{process\_apdu()} implementation itself.
\end{itemize}

Only after these checks does the kernel update its single selected
context slot. This is the precise meaning of selection in the current
implementation: \texttt{SELECT} does not install the application, does
not invoke the \texttt{install()} callback, and does not yet grant a
general execution mode. It merely resolves an AID to an embedded
administrative application, loads that application, and records it as
the selected execution context for subsequent command dispatch.

The later command path then uses kernel-side trampolines to enter the
selected application. These trampolines load the application's
global-pointer context before the callback and restore the kernel state
after the callback returns. This context-switching rule is the critical
mechanism that allows the kernel and the selected FAE image to coexist
without treating the application as part of the kernel's own relocation
domain.

\subsection{Rustlet Runtime, ABI and Isolation}
\label{sec:impl-rustlet-runtime}
% --- Rustlet runtime and isolation implementation ---
% This subsection should describe the concrete implementation of the
% abstract Rustlet model introduced in Chapter 7.
% It should cover:
% - the current shared ABI based on RustletCtx;
% - the declare_rustlet! macro and its generated start entry point;
% - the selected-application descriptor and vtable;
% - the current core::isolation session model;
% - MPU planning and region programming;
% - the concrete gate region layout;
% - the ARM M-profile SVC / MemManage / resume path;
% - how install, process_apdu, exit, panic and faults converge;
% - the current limits of state persistence and serialization.

The Rustlet runtime is split between two binaries. The kernel owns selection,
registry lookup, APDU transport, MPU programming and syscall dispatch. The
Rustlet image owns the application state, the Rustlet heap, the typed APDU
framework and the exported entry points expected by the kernel. The boundary
between the two is deliberately a C-style ABI rather than a native Rust trait
object, because the kernel and the Rustlet FAE image have distinct relocation
contexts.

The first entry point of a Rustlet image is the FAE entry at offset
\texttt{0}. The kernel enters this image with the shared
\texttt{RustletCtx} pointer and the application global-pointer context. The
Rustlet runtime initializes its local runtime globals, aligns the private heap
storage declared by the macro, builds a \texttt{SelectedAppDescriptor}, and
returns that descriptor through the dedicated runtime-return syscall. The
descriptor contains:
\begin{itemize}
\item an opaque runtime state pointer;
\item a selected-application vtable containing \texttt{install()} and
  \texttt{process\_apdu()} handlers;
\item an optional Security-Domain vtable for Rustlet-backed Security Domains;
\item the Rustlet heap window exported to the kernel.
\end{itemize}

The \texttt{declare\_rustlet!} macro is the developer-facing source of this
low-level structure. It generates the exported \texttt{start()} function,
allocates the static Rustlet heap storage, adapts the user installation
function into a boxed runtime instance, and wraps that instance in the runtime
persistence adapter. \texttt{declare\_app!} is only a compatibility alias.
\texttt{declare\_security\_domain!} follows the same pattern, but also exports
the \texttt{sddispatch} vtable used by
\texttt{RustletSecurityDomainProxy}.

All APDU exchange uses one shared \texttt{RustletCtx}. Before a Rustlet call,
the kernel fills the APDU header fields and stages any incoming payload in the
context payload buffer. During the call, Rustlet code normally uses the typed
\texttt{Apdu<State>} wrapper described in Chapter~7, but the ABI object remains
the same shared buffer. If the Rustlet switches to outgoing mode, the same
payload area becomes the response area. On return, the kernel publishes the
outgoing length and status word in its transport session. In firmware, the
transport payload and the Rustlet payload are two views of the same gate-region
storage, so this publication does not normally copy the payload. Normal command
completion clears the payload half. An abrupt Rustlet exit clears the complete
shared page, and every later Rustlet entry reinitializes the control half before
user code can observe it.

The shared page has two implementation-level invariants that are deliberately
stronger than the public Rustlet API. First, the APDU payload half is the only
place where Rustlet command or response payload bytes are staged. Protocol
layers must not introduce larger stack-local APDU mirrors on the hot path.
Second, the control half is also the kernel's secondary scratch buffer before a
Rustlet entry. Secure-channel code may use it while parsing, authenticating or
wrapping APDUs, but every Rustlet entry path must reset it to the ABI layout
before user code can observe the page. This is especially important when the
scratch bytes contain decrypted data, MAC material or serialized state.

The application call path is normalized around two vtable handlers. The
\texttt{install()} handler constructs the resident Rustlet instance from the
\texttt{INSTALL [for install]} parameters exposed by the kernel. The
\texttt{process\_apdu()} handler loads the current serialized state if
necessary, calls the Rustlet implementation, and saves the resident state again
before returning. A Security-Domain Rustlet uses the same selected-application
handlers for ordinary APDUs and the separate \texttt{sddispatch} handler for
kernel-initiated management, secure-channel and secure-messaging hooks.

Isolation is planned before the kernel enters the Rustlet. The current MPU plan
contains the Rustlet text window, one writable RAM block containing stack,
relocated data and heap, and the gate region. The writable Rustlet RAM block is
mapped read/write and non-executable. The Rustlet text window is mapped
executable only while the CPU is actually executing Rustlet code. When the
Rustlet traps into the kernel through a syscall, the kernel switches to a
``kernel serving Rustlet'' phase and disables execution from the active Rustlet
text regions before servicing the request. Execution permission is restored
only when control is returned to the Rustlet.

The gate region contains the shared ABI buffer and the small target-specific
entry/resume stubs required to cross the privileged/unprivileged boundary. It
is the only deliberately shared execution bridge. The rest of the Rustlet
memory map is explicit: Rustlet code cannot read or write kernel memory, cannot
execute from its writable RAM block, and cannot execute from its own text while
the kernel is servicing a syscall.

Rustlet termination paths converge through the same isolation machinery. A
normal handler return, an explicit Rustlet exit, a panic, and a recoverable MPU
fault all redirect execution back to the kernel resume point with a small
return kind and status value. For ARMv7-M targets this is primarily driven by
the MemManage path; for ARMv6-M targets without a separate MemManage exception,
the board port may translate the relevant HardFault cases into the same
recoverable application-fault path. In both cases, a fault while no Rustlet call
is active remains a kernel fault, not an application status word.

The current persistence boundary is intentionally at the serialized Rustlet
state level. The runtime adapters can load and save the resident state through
the shared context, and the kernel stores that serialized state in the registry
object for the selected instance. This is not yet a full non-volatile storage
backend: without persistent registry storage, rebooting the firmware rebuilds
the initial predeployment state. The important invariant is nevertheless
already in place: Rustlet execution sees an instance state, not a global static
singleton, and the registry object is the place where that state will become
durable once flash-backed persistence is implemented.

\section{APDU Parsing and Command Dispatch}
\label{sec:impl-apdu-dispatch}
% --- APDU parsing and command dispatch ---
% This section describes how APDU commands are decoded and dispatched
% to the appropriate handlers (LOAD, INSTALL, DELETE, etc.)

The transport mechanics of the \texttt{T=0} exchange are described in
Section~\ref{sec:impl-apdu-t0}. The APDU-facing implementation surfaces
are summarized in Section~\ref{sec:impl-apdu-surfaces}. Once the command
header has been captured, the dispatcher receives an
\texttt{ApduCommand} through the \texttt{ApduLayer} boundary, derives an
ephemeral kernel-side \texttt{Apdu} view from it, and then applies the
dispatch rules. This keeps the transport representation internal to the
base \texttt{T=0} manager and reserves the space for higher protocol
layers such as SCP03.

The dispatch layer is deliberately small and distinguishes four cases before
ordinary application-command routing:
\begin{itemize}
\item \texttt{GET RESPONSE} remains entirely transport-owned;
\item \texttt{SELECT} remains kernel-owned and updates the selected
  context as described in Section~\ref{sec:impl-selected-context};
\item GlobalPlatform management commands that mutate or query platform state
  remain kernel-owned and are authorized by the active Security Domain;
\item every other command is forwarded to the currently selected application
  through the selected-application ABI.
\end{itemize}

This organization keeps the boundary between transport, protocol-layer
composition, and command logic explicit. The base \texttt{T=0} manager
remains responsible for byte acquisition, deferred \texttt{GET RESPONSE}
handling, wrong-\texttt{Le} handling through \texttt{6Cxx}, and final
\texttt{SW1/SW2} emission. The secure-channel layer is responsible for
establishment, protected-command unwrapping, and protected-response wrapping.
The management dispatcher is responsible for the kernel-owned GlobalPlatform
operations. Only after those layers have declined a command does the ordinary
selected-application ABI receive it.

The current ABI-level exchange object is a fixed-layout
\texttt{RustletCtx} structure containing exactly two \texttt{256}-byte halves:
\begin{itemize}
\item a control/secondary buffer containing the ABI version, APDU
  header/status bytes, flags, the one-byte APDU data length and serialized
  state;
\item one fixed \texttt{256}-byte APDU payload buffer.
\end{itemize}

If the APDU requires command data reception, the kernel receives that
payload and stages it into the incoming view of this shared buffer. If
the Rustlet switches the buffer to outgoing mode, the same payload area
is then interpreted as the response buffer.

The APDU payload buffer is physically \texttt{256} bytes long, but the current
short-APDU ABI exposes at most \texttt{255} useful bytes. This is intentional:
the length field is one byte, and extended APDU support must be introduced as
an ABI evolution rather than by silently overloading the current layout. The
serialized-state slice lives in the control/secondary half and is currently
limited to \texttt{244} bytes after ABI metadata.

The kernel-to-Rustlet bridge is normalized around the
\texttt{SEApdu} trait. In practice, three APDU-facing objects implement
that common interface:
\begin{itemize}
\item the kernel-side \texttt{Apdu} wrapper;
\item the Rustlet-side \texttt{RustletCtx};
\item the kernel bridge object that couples the transport-owned APDU
  state with the shared Rustlet ABI buffer while a Rustlet call is in
  progress.
\end{itemize}

APDU command logic manipulates the current \emph{card-side APDU
session} through the common APDU interface, while transport remains
kernel-owned.

Dispatch therefore follows this kernel-owned rule:
\begin{itemize}
\item \texttt{GET RESPONSE} stays transport-owned;
\item \texttt{SELECT} stays kernel-owned and resolves an instance AID in
  the registry;
\item secure-channel establishment and protected-message framing stay in
  \texttt{SecureChannelLayer};
\item management commands such as \texttt{INSTALL}, \texttt{LOAD},
  \texttt{DELETE}, \texttt{PUT KEY}, \texttt{GET DATA}, and
  \texttt{GET STATUS}, \texttt{STORE DATA}, and \texttt{SET STATUS} stay
  kernel-owned and are authorized by the active Security Domain;
\item otherwise, the kernel calls the selected application's
  \texttt{process\_apdu()} entry point.
\end{itemize}

For \texttt{INSTALL}, the kernel parses the
GlobalPlatform \texttt{INSTALL [for install]} body as a sequence of
length-value fields:
\begin{itemize}
\item Executable Load File AID (package AID);
\item Executable Module AID (applet AID);
\item Application Instance AID;
\item privileges;
\item install parameters.
\end{itemize}

The kernel then resolves the target Rustlet from the package/applet AID
pair, updates the current instance AID, and only exposes the
\emph{install parameters} to the Rustlet-side \texttt{install()}
callback. This means that the Rustlet constructor sees the installation
payload relevant to the application itself, while the GlobalPlatform
deployment envelope remains kernel-owned.

This rule is still intentionally narrow. The goal is to validate the
boundary between the kernel and the selected administrative application
before encoding the full GlobalPlatform routing matrix. In particular,
\texttt{SELECT} itself never invokes \texttt{install()}; \texttt{INSTALL}
is an explicit command interpreted according to the GlobalPlatform body
structure before the
application callback is entered.

In the current Rustlet model, \texttt{install()} is the application-side
constructor of the selected context. The ABI has no separate definition
object.

When the selected application returns, the kernel copies outgoing bytes only
when the bridge is backed by distinct storage. The normal firmware path simply
publishes the outgoing length for the already shared payload and then resumes
the normal \texttt{T=0} completion logic. This means that deferred
\texttt{GET RESPONSE} ownership, procedure-byte emission, and
\texttt{6Cxx}/\texttt{61xx} completion rules remain entirely in the
kernel while the selected application supplies the command semantics.

This is the implementation form of the doctrine described in
Chapter~5. Oxide SE does not currently invert the hook by making the selected
Security Domain parse the whole GlobalPlatform command grammar and then call
registry services. The kernel parses the management envelope, preserves the
registry invariants, and asks the active Security Domain to authorize the
operation in its name.

The following constraints remain explicit:
\begin{itemize}
\item there is only one effective logical channel and one selected context
  slot;
\item dynamically loaded packages are not yet implemented; packages are
  currently declared by the predeployment manifest and embedded into the
  firmware image;
\item the incoming-length heuristic is still a short-APDU simplification
  and does not yet encode the full distinction between command cases;
\item the future persistent registry backend still has to provide
  non-volatile atomicity for the registry-object operations described above.
\end{itemize}

\section{Security Domains and Access Control}
\label{sec:impl-security-domain}
% --- Privilege and security domain enforcement ---
% This section describes:
% - privilege checking
% - Security Domain authority
% - access control enforcement
% - delegation model

Security-Domain authority is represented by registry objects whose
\texttt{kind} is \texttt{SecurityDomain}. The concrete behavior behind such an
object is selected by its backend:
\begin{itemize}
\item \texttt{NullSecurityDomain} is a permissive development backend used for
  bring-up and dynamic Rustlet experiments without secure-channel keys;
\item \texttt{KernelSecurityDomain} keeps Security-Domain policy and SCP state
  inside the firmware;
\item \texttt{RustletSecurityDomainProxy} keeps the kernel-side authority
  checks and delegates backend-specific policy or SCP processing to an isolated
  Rustlet Security Domain.
\end{itemize}

All three backends share the same kernel-owned base authority. Their
three-byte privilege field is interpreted as GlobalPlatform administrative
authority over operation families: load, install, delete, lifecycle
management, key management, delegated management, token or receipt handling,
card-level lock/termination operations, and related management functions as
they become implemented. The backend may refuse an operation even when the base
bit allows it, but it may not turn a missing base privilege into an allowed
kernel mutation.

This gives Rustlet-backed Security Domains their intended shape. The Rustlet
can implement richer local policy and cryptographic behavior, and it may store
additional serialized state of its own. The proxy nevertheless enforces the
kernel-owned privilege floor first. If the recorded three-byte privilege set
does not permit an operation family, the kernel refuses the operation without
asking the Rustlet backend to override that decision.

Delegation is checked at Security-Domain creation time. A supplementary
Security Domain is installed under the currently active administrative
Security Domain and receives that active domain as \texttt{parent\_sd\_aid}.
The requested privilege bytes are accepted only if they are no broader than the
authority that the parent may delegate. The resulting object then participates
in the same filtered registry view as every other managed object.

In other words, Security Domains are first-class registry objects with
selectable AIDs and backend-specific behavior, but the registry, privilege
non-escalation checks, and platform-state mutations remain kernel-mediated.

\section{Secure Channel Command Handling}
\label{sec:impl-scp-commands}
% --- Secure channel command handling ---
% This section describes:
% - handling of INITIALIZE UPDATE
% - handling of EXTERNAL AUTHENTICATE
% - secure messaging extensions linked to the secure channel state

The protocol shape of SCP03 and SCP11 is defined in
Section~\ref{sec:protocol-scp}. The longer examples, including the
Amendment~D command and response framing and the
representative SCP03 and SCP11a/b/c APDU flows, are kept in
Appendix~\ref{app:scp-data-objects}. The present section is narrower: it
describes where those commands enter the reference kernel and how they are
connected to Security-Domain policy, registry-backed key material, and APDU
dispatch.

\subsection{Establishment Dispatch}
\label{sec:impl-scp-establishment-dispatch}

Secure-channel command handling is centralized in
\texttt{SecureChannelLayer}. This layer wraps the base \texttt{T=0} APDU
manager. It receives complete APDU commands from the transport, handles
secure-channel establishment commands before ordinary dispatch, unwraps
protected commands when the secure-messaging bit is set in the CLA byte, and
wraps protected responses before the transport emits the final response APDU.
The layer is intentionally profile-neutral: it knows which APDU opcodes belong
to secure-channel establishment, but it delegates the cryptographic decision to
the active Security Domain backend.

The establishment candidates are:
\begin{itemize}
\item \texttt{INITIALIZE UPDATE} and \texttt{EXTERNAL AUTHENTICATE} for SCP03;
\item \texttt{PERFORM SECURITY OPERATION} and \texttt{MUTUAL AUTHENTICATE} for
  SCP11a and SCP11c;
\item \texttt{INTERNAL AUTHENTICATE} for SCP11b.
\end{itemize}

The set of accepted candidates is controlled by the effective build
configuration described in Section~\ref{sec:impl-scp-profiles}. If SCP11a and
SCP11b are compiled and enabled but SCP11c is not, the layer rejects an SCP11c
establishment attempt before it can become an authenticated session. If no SCP
protocol is enabled, the layer becomes a pass-through boundary except for the
identity/test mode used by kernel validation.

Once an establishment APDU has been recognized, the layer consumes its incoming
payload, builds a \texttt{SecureChannelEstablishmentCommand}, and calls the
current Security Domain. Kernel-native Security Domains resolve static keys or
SCP11 material inside the firmware. Rustlet-backed Security Domains receive the
same request through the proxy path and may perform the profile-specific
cryptographic work in user-land. In both cases, the active Security Domain is
the authority that owns the session state; \texttt{SecureChannelLayer} only
owns the APDU boundary.

\subsection{SCP03 INITIALIZE UPDATE Handling}
\label{sec:impl-scp-init}
% --- INITIALIZE UPDATE handling ---
% This subsection describes:
% - host challenge processing
% - card challenge generation
% - key version selection
% - session key derivation (SCP02 / SCP03)
% - construction of response cryptogram
% - preparation of secure session context

\texttt{INITIALIZE UPDATE} is the first SCP03 establishment command. The
reference APDU shape and the expected protocol fields are described in
Section~\ref{sec:protocol-scp}; the exact SCP03 KDF material is detailed in
Appendix~\ref{app:scp-data-objects}. On the implementation path, the secure
channel layer recognizes \texttt{INS = 50} on a GlobalPlatform secure-channel
CLA and forwards the command body to the active Security Domain.

The Security Domain then performs the SCP03-specific work:
\begin{itemize}
\item parse the key version, key identifier and host challenge carried by the
  command;
\item resolve the corresponding static key objects in the registry, scoped by
  the Security Domain instance;
\item generate the card challenge and derive the session keys;
\item compute the card cryptogram and build the response payload;
\item store the initialized but not yet authenticated session state.
\end{itemize}

This split keeps static-key storage out of the protocol core. The pure SCP03
engine receives explicit key material and challenges; it does not know whether
those keys came from the root Security Domain, a supplementary kernel-native
Security Domain, or a Rustlet-backed Security Domain proxy. Conversely, the
registry and Security-Domain layers do not implement CMAC/KDF internals
themselves. They resolve authority and key material, then call the pure engine.

At the end of a successful \texttt{INITIALIZE UPDATE}, the secure channel is
not yet open. The session material is staged, but protected APDUs are still
rejected until \texttt{EXTERNAL AUTHENTICATE} succeeds. This distinction is
important for replay and stale-session handling: a second establishment
sequence replaces the pending material for that Security Domain and logical
channel rather than silently authenticating an older chain.

\subsection{SCP03 EXTERNAL AUTHENTICATE Handling}
\label{sec:impl-scp-auth}
% --- EXTERNAL AUTHENTICATE handling ---
% This subsection describes:
% - host cryptogram verification
% - security level activation
% - transition to an authenticated secure-channel state
% - failure handling and session rejection

\texttt{EXTERNAL AUTHENTICATE} completes SCP03 establishment. The command
contains the requested security level, the host cryptogram computed over the
session material produced by \texttt{INITIALIZE UPDATE}, and the transmitted
initial C-MAC required by SCP03 Amendment~D. The secure-channel layer recognizes
\texttt{INS = 82} with the SCP03 establishment shape and forwards the request
to the active Security Domain.

The Security Domain verifies both the host cryptogram and the C-MAC against the
pending SCP03 session. The C-MAC is computed from an all-zero chaining value
over the \texttt{84 82 P1 00 Lc} command header and the host cryptogram; the
\texttt{Lc} byte includes the transmitted C-MAC. If verification succeeds, the
Security Domain records the authenticated security level and initializes the
secure-messaging state: command MAC chain, response MAC chain, encryption
counters, and profile-specific MAC length. The full sixteen-byte
\texttt{EXTERNAL AUTHENTICATE} C-MAC seeds both MAC chains, including for S8
where only eight bytes are transmitted. If verification fails, the pending
session remains unauthenticated and the command returns a failure status rather
than leaving a partially authenticated session behind.

The security level selected here controls later protected APDUs. In the
Oxide SE SCP03 profile, C-MAC authenticates commands and C-ENC decrypts command
data; responses remain plain because response-protection levels are explicitly
profiled out. The APDU dispatcher does not duplicate this matrix. It asks the
active Security Domain to unwrap commands and consults the authenticated
session state before deciding whether response wrapping applies. The exact
selected values and rejections are stated in
Section~\ref{sec:impl-scp-selected-levels}.

\subsection{SCP11 PSO, MUTUAL AUTHENTICATE and INTERNAL AUTHENTICATE Handling}
\label{sec:impl-scp11-establishment}

SCP11a, SCP11b and SCP11c enter the implementation through the common
secure-channel establishment dispatcher. Unlike SCP03, they do not use the
pair formed by \texttt{INITIALIZE UPDATE} and
\texttt{EXTERNAL AUTHENTICATE}. Their normative command sequences are described
in Section~\ref{sec:protocol-scp} and illustrated in
Appendix~\ref{app:scp-data-objects}. The implementation section should
therefore be read as a routing and state-management description, not as a
second copy of the SCP11 standard text.

For SCP11a and SCP11c, \texttt{PERFORM SECURITY OPERATION} is handled as a
profile-specific establishment command. The secure-channel layer recognizes
\texttt{INS = 2A} when SCP11 support is enabled, receives the incoming BER-TLV
certificate or key material, and forwards it to the active Security Domain. The
Security Domain validates the profile-specific material and records the
intermediate SCP11 establishment state required by the following
\texttt{MUTUAL AUTHENTICATE}. A repeated or stale PSO sequence replaces or
rejects that intermediate state according to the Security Domain policy; it
must not silently authenticate an older receipt chain.

\texttt{MUTUAL AUTHENTICATE} then opens SCP11a or SCP11c. The secure-channel
layer recognizes the non-SCP03 \texttt{INS = 82} shape and forwards the APDU
body to the same active Security Domain. The backend parses the CRT and public
key material, computes the profile-specific ECDH contribution, builds the
GlobalPlatform SharedInfo value, derives the secure-messaging keys, verifies
the expected receipt or cryptogram, and records the authenticated session
state. Kernel-native and Rustlet-backed Security Domains share the same APDU
boundary; what differs is where the profile-specific cryptography and policy
are executed.

SCP11b opens through \texttt{INTERNAL AUTHENTICATE}. The secure-channel layer
recognizes \texttt{INS = 88} only when SCP11b is enabled, then delegates the
body to the active Security Domain. The backend performs the SCP11b
card-authentication key agreement, derives the same kind of secure-messaging
state used by the other profiles, and records the session as authenticated.
Because SCP11b does not authenticate the OCE inside the SCP11 exchange itself,
any additional authorization policy belongs to the selected Security Domain
rather than to \texttt{SecureChannelLayer}.

After successful SCP11 establishment, the command path converges with SCP03:
later protected APDUs use the same raw payload-and-MAC framing, and the
dispatcher remains profile-neutral after unwrapping.

\subsection{Protected APDU Unwrap and Wrap}
\label{sec:impl-scp-ext}

The protected-APDU path is the implementation counterpart of the
secure-messaging grammar introduced in Section~\ref{sec:protocol-scp} and
specified byte-for-byte in Appendix~\ref{app:scp-data-objects}. Its purpose is
not to implement application commands. It is the boundary that turns a
wire-level protected APDU into the clear command semantics consumed by the
kernel dispatcher, then turns the logical response back into a protected
GlobalPlatform response body.

\subsubsection{Command classification}

Protected command handling starts from the CLA byte. In the command families
accepted by the reference kernel, bit \texttt{0x04} marks a
secure-messaging command. If this bit is present while no authenticated session
is open for the active Security Domain, the command is rejected with a
conditions-not-satisfied status before it reaches management or Rustlet
dispatch.

For SCP03, a command without that bit continues through the clear-command path
and is interpreted under the authority policy described in
Section~\ref{sec:impl-apdu-dispatch}. SCP11 applies the stricter Amendment~F
session rule: while secure messaging is active, clear \texttt{SELECT}
terminates the session and is then processed normally; every other clear
command aborts the session and returns \texttt{6985}. A protected management
command that unwraps successfully is then checked against the peer/profile
matrix in Section~\ref{sec:impl-scp-selected-levels} before any registry
mutation or Security Domain hook.

The identity secure-channel mode used by kernel validation is a special build
profile. It leaves the APDU body untouched so that low-level transport and
dispatcher tests can exercise the same layer ordering without cryptographic
state. Production SCP03 and SCP11 profiles do not use that bypass.

\subsubsection{Protected command unwrapping}

For a protected command, \texttt{SecureChannelLayer} first receives the
incoming APDU bytes into the shared APDU buffer. As described in
Appendix~\ref{app:scp-data-objects}, it separates the trailing
profile-dependent C-MAC from the preceding clear or encrypted command data.
The APDU \texttt{Le} remains in the APDU encoding. The layer rejects missing
MAC bytes and oversized protected fields before the command can be dispatched.

The layer then reconstructs the authenticated command string from the APDU
header and the authenticated prefix of the protected body. This is the byte
string whose MAC-chain semantics are defined by the selected SCP profile. The
actual verification, counter update, decryption, and padding removal are
delegated to the active Security Domain through its
\texttt{unwrap\_command()} hook. This keeps profile-specific cryptography,
registry-backed keys, and Security-Domain policy out of the transport layer.

If unwrapping succeeds, the protected incoming payload in the APDU buffer is
replaced by the clear payload returned by the Security Domain. From that point
onward, management handlers and Rustlet dispatch see the same logical command
they would have seen in clear. They do not need to know whether the external
wire representation contained clear or encrypted protected data and its
trailing MAC; secure messaging remains a boundary concern.

The unwrap path must preserve the shared-buffer discipline. The protected APDU
body is parsed from the APDU payload buffer, while temporary clear/protected
material should use the secondary half of \texttt{RustletCtx} or another
explicit kernel-owned scratch area, not unbounded stack arrays. Once the
logical command is reconstructed, the APDU payload buffer contains only the
clear command payload that will be seen by management handlers or by the
selected Rustlet.

\subsubsection{Protected response wrapping}

Response handling mirrors the command path. If the command was protected, the
logical response bytes and the logical status word are given to the active
Security Domain through its \texttt{wrap\_response()} hook. The Security Domain
applies the authenticated session state and returns a protected response body:
clear or encrypted response data followed by the R-MAC when the selected
security level requires one. The logical status remains the actual trailing
APDU status word and is included in the R-MAC calculation.

Whether wrapping is required is captured from the authenticated session while
the command is being unwrapped, before ordinary dispatch. The resulting policy
bit is stored in the current \texttt{ApduCommand} and consumed during response
completion; it cannot leak into the next command because it has the same
lifetime as that APDU. This ordering is significant for a Rustlet-backed
Security Domain: querying its state uses \texttt{sddispatch}, which owns the
same shared APDU page and would otherwise overwrite an application response
that has already been produced.

The transport emits the logical application or management status as the
actual trailing status word. When R-MAC is active, that status is
authenticated together with the protected response data, as illustrated in
Appendix~\ref{app:scp-data-objects}.

\subsubsection{Profile neutrality and replay invariants}

The unwrap/wrap path is common to SCP03 and SCP11 sessions. What differs
between profiles is the establishment sequence and the session material from
which MAC and encryption keys were derived. After establishment, the command
dispatcher deliberately remains profile-neutral: it handles a clear command
view and relies on the active Security Domain to enforce the selected SCP
profile, security level, MAC length, counters, receipt state, and replay
policy.

The invariant is simple: a protected APDU that fails MAC verification, uses a
stale chain value, exceeds the bounded protected-payload limits, or targets a session
that is not authenticated must fail before ordinary command logic sees it.
Only successfully unwrapped commands may mutate registry state, call Rustlet
code, or produce protected responses.

The implementation exposes one status-word contract for both the native and
Rustlet-backed Security Domain paths:
\begin{description}
\item[\texttt{6700}] a profile-dependent length is invalid or a bounded output
      cannot contain the requested transformation;
\item[\texttt{6A80}] an SCP11 CRT or another command data structure is
      malformed;
\item[\texttt{6A86}] \texttt{P1/P2} is invalid or selects an option excluded
      by the compiled profile;
\item[\texttt{6A88}] the referenced SCP03 keyset or SCP11 ECKA/CA key does not
      exist for the active Security Domain;
\item[\texttt{6600}] SCP11 certificate verification failed;
\item[\texttt{6300}] SCP03 host authentication failed;
\item[\texttt{6982}] secure-messaging authentication failed, including a stale
      MAC chain, stale receipt, or replay;
\item[\texttt{6985}] the command is valid but its protocol precondition is not
      satisfied;
\item[\texttt{6D00}] the establishment instruction is outside the compiled
      SCP profile.
\end{description}

MAC chains and encryption counters are transactional. The engines derive a
candidate next state, complete every bounds check and cryptographic transform,
and only then publish that state. A rejected protected APDU therefore cannot
consume one chain value or counter. Structural errors leave an authenticated
session available for a corrected APDU. A cryptographic secure-messaging
failure returns \texttt{6982}, erases the session, and causes later protected
commands to return \texttt{6985} until establishment is completed again.
Likewise, a failed \texttt{EXTERNAL AUTHENTICATE} cannot be retried without a
new \texttt{INITIALIZE UPDATE}.
